
\section*{Round 2 continuation for Erd\H{o}s problems \#319, \#359, \#361}

This file is written as a continuation of the Round-1 write-up.  I explicitly indicate below whenever I use a Round-1 result as a black box.

\bigskip
\hrule
\bigskip

\section*{Problem \#361}

\subsection*{1) ROUND-2 OBJECTIVE}
\textbf{Path (C): obstruction/correction.}  Round 1 completely solved the case $c\ge 1$, but left the whole range $0<c<1$ open.  The most promising Round-2 direction is therefore to give genuinely new \emph{rigorous constructions} for $c<1$, especially constructions that explain why arithmetic properties of $n$ should matter.

I will prove two new lower-bound mechanisms:

\begin{itemize}
\item an interval construction valid for all $c<1$;
\item a new \emph{hybrid modular/large-element construction} for $1/2<c<1$.
\end{itemize}

The hybrid construction depends explicitly on the smallest integer $q>1$ not dividing $n$, and this gives a rigorous arithmetic obstruction that was absent from Round 1.

\subsection*{2) ROUND-1 FOUNDATION USED}
I use one Round-1 theorem as a black box.

\begin{itemize}
\item Round 1 proved that if $m\ge n$, then the maximum size of $A\subseteq \{1,\dots,m\}$ with $n\notin \Sigma(A)$ is exactly
\[
m-\left\lceil \frac n2\right\rceil.
\]
In particular, for fixed $c\ge 1$ and $m=\lfloor cn\rfloor$,
\[
G_c(n)=\lfloor cn\rfloor-\left\lceil \frac n2\right\rceil
= \left(c-\frac12\right)n+O(1).
\]
\end{itemize}

\subsection*{3) NEW INSIGHT / TOOL (ROUND-2)}
The new ingredient is that, when $m<n$, two different obstructions can be superimposed:

\begin{itemize}
\item \textbf{size obstruction:} if all chosen elements lie in a short interval $(n/b,m]$, then any $(b-1)$-term sum is too small and any $b$-term sum is too large;
\item \textbf{modular obstruction plus one-large-element principle:} if all elements above $n/2$ avoid the residue class $n \pmod q$, while all elements at most $n/2$ are multiples of $q$, then any subset summing to $n$ is impossible.
\end{itemize}

The second point is genuinely new in this write-up and produces a clean arithmetic lower bound depending on the smallest $q>1$ with $q\nmid n$.

\subsection*{4) ATTACK PLAN (ROUND-2)}
The exact Round-1 gap was the complete absence of rigorous theorems for $c<1$.

I will now prove:

\begin{itemize}
\item[(i)] a general interval-construction lower bound for $m<n$;
\item[(ii)] a new hybrid lower bound for $m>n/2$;
\item[(iii)] the asymptotic form of the hybrid lower bound
\[
G_c(n)\ge \left(c-\frac12+\frac{1-c}{q(n)}\right)n+O(1),
\]
where $q(n)$ is the smallest integer $>1$ not dividing $n$.
\end{itemize}

This does not settle the problem for $c<1$, but it gives a sharply formulated obstruction: below $c=1$, the extremal behaviour already depends on the arithmetic of $n$ at the level of explicit linear constructions.

\subsection*{5) WORK (ROUND-2)}

For this section write
\[
G(n,m):=\max\{|A|:\ A\subseteq \{1,\dots,m\},\ n\notin \Sigma(A)\}.
\]
Thus $G_c(n)=G(n,\lfloor cn\rfloor)$.

\paragraph{Proposition 361.1 (interval construction).}
Let $b\ge 2$ be an integer and suppose
\[
m<\frac{n}{b-1}.
\]
Define
\[
A_b:=\left\{\left\lfloor \frac nb\right\rfloor+1,\ \left\lfloor \frac nb\right\rfloor+2,\ \dots,\ m\right\}.
\]
Then $n\notin \Sigma(A_b)$.  Consequently,
\[
G(n,m)\ge \max\!\left(0,\ m-\left\lfloor \frac nb\right\rfloor\right).
\]

\paragraph{Proof.}
Every element of $A_b$ is strictly larger than $n/b$.  Therefore any sum of $b$ elements of $A_b$ is strictly larger than $n$.

On the other hand, every element of $A_b$ is at most $m$, and by assumption
\[
(b-1)m<n.
\]
So any sum of at most $b-1$ elements of $A_b$ is strictly smaller than $n$.

Hence no subset of $A_b$ can sum to $n$.  The cardinality statement is immediate.
\hfill $\square$

\paragraph{Remark.}
If $1/b\le c<1/(b-1)$ and $m=\lfloor cn\rfloor$, Proposition 361.1 gives
\[
G_c(n)\ge \left(c-\frac1b\right)n+O(1).
\]

\paragraph{Proposition 361.2 (hybrid modular/large-element construction).}
Assume
\[
\frac n2 < m < n.
\]
Let $q>1$ be any integer with $q\nmid n$, and define
\[
A_q(n,m):=
\{a\in \{1,\dots,\lfloor n/2\rfloor\}: q\mid a\}
\ \cup\
\{a\in \{\lfloor n/2\rfloor+1,\dots,m\}: a\not\equiv n \pmod q\}.
\]
Then
\[
n\notin \Sigma(A_q(n,m)).
\]

\paragraph{Proof.}
Let $S\subseteq A_q(n,m)$ be any subset.

\smallskip
\noindent
\emph{Case 1: $S$ contains no element greater than $n/2$.}
Then every element of $S$ is divisible by $q$, so
\[
\sum_{a\in S} a \equiv 0 \pmod q.
\]
Since $q\nmid n$, we have $n\not\equiv 0\pmod q$, so $\sum_{a\in S}\neq n$.

\smallskip
\noindent
\emph{Case 2: $S$ contains at least two elements greater than $n/2$.}
Then
\[
\sum_{a\in S} a > \frac n2 + \frac n2 = n,
\]
so again $\sum_{a\in S}\neq n$.

\smallskip
\noindent
\emph{Case 3: $S$ contains exactly one element $y>n/2$.}
Every other element of $S$ is divisible by $q$, so
\[
\sum_{a\in S} a \equiv y \pmod q.
\]
But by construction $y\not\equiv n\pmod q$.  Hence $\sum_{a\in S}\neq n$.

These three cases exhaust all possibilities, so indeed $n\notin \Sigma(A_q(n,m))$.
\hfill $\square$

\paragraph{Corollary 361.3.}
Let $q(n)$ denote the smallest integer $>1$ that does not divide $n$.  If $1/2<c<1$ and $m=\lfloor cn\rfloor$, then
\[
G_c(n)\ge \left(c-\frac12+\frac{1-c}{q(n)}\right)n+O(1).
\]

\paragraph{Proof.}
Apply Proposition 361.2 with $q=q(n)$.  The size of the constructed set is
\[
|A_{q(n)}(n,m)|
=
\#\{a\le \lfloor n/2\rfloor:\ q(n)\mid a\}
+
\#\{\lfloor n/2\rfloor<a\le m:\ a\not\equiv n\pmod{q(n)}\}.
\]
The first term is
\[
\frac{n}{2q(n)}+O(1),
\]
and the second term is
\[
\left(m-\frac n2\right)\left(1-\frac1{q(n)}\right)+O(1)
=
\left(c-\frac12\right)\left(1-\frac1{q(n)}\right)n+O(1).
\]
Adding,
\[
|A_{q(n)}(n,m)|
=
\left(c-\frac12+\frac{1-c}{q(n)}\right)n+O(1).
\]
Therefore
\[
G_c(n)\ge \left(c-\frac12+\frac{1-c}{q(n)}\right)n+O(1).
\]
\hfill $\square$

\paragraph{Examples when $c=3/4$.}
The new lower bound gives the following explicit infinite families.

\begin{itemize}
\item If $n$ is odd, then $q(n)=2$, so
\[
G_{3/4}(n)\ge \frac{3}{8}n+O(1).
\]
This is exactly the ``all even numbers'' construction.

\item If $n$ is even and $3\nmid n$, then $q(n)=3$, so
\[
G_{3/4}(n)\ge \frac13 n+O(1).
\]

\item If $6\mid n$ but $4\nmid n$, then $q(n)=4$, so
\[
G_{3/4}(n)\ge \frac{5}{16}n+O(1).
\]
\end{itemize}

Thus, already at the level of clean explicit constructions, the lower-bound constant changes with the divisibility pattern of $n$.

\subsection*{6) ADVERSARIAL VERIFICATION}
I check both new constructions carefully.

\begin{itemize}
\item \textbf{Interval construction:}  
The strict inequalities are essential.  I used $m<n/(b-1)$ to force every $(b-1)$-term sum to be $<n$, and I used the lower endpoint $\lfloor n/b\rfloor+1$ to force every $b$-term sum to be $>n$.  No boundary case has been glossed over.

\item \textbf{Hybrid construction:}  
The only subtle point is possible overlap between the ``lower multiples of $q$'' and the upper interval.  I avoided this by defining the two pieces on disjoint ranges:
\[
[1,\lfloor n/2\rfloor]
\quad\text{and}\quad
[\lfloor n/2\rfloor+1,m].
\]
So the case distinction by number of upper elements is exact.

\item \textbf{Could one upper element plus lower multiples of $q$ still sum to $n$?}  
No.  Modulo $q$, such a sum would be congruent to the upper element itself, and by construction every upper element avoids the residue class $n \pmod q$.

\item \textbf{Does the new lower bound really interpolate correctly between known constructions?}  
Yes.  When $q(n)=2$ (for odd $n$), the hybrid set is exactly the set of even numbers, giving the parity construction.  When $q(n)$ is large, the gain over the pure interval construction becomes small, as expected.
\end{itemize}

\subsection*{7) FINAL}
\textbf{UNRESOLVED (BUT STRICTLY ADVANCED).}

Round 2 adds genuine new mathematics below the previously solved threshold $c=1$:

\begin{itemize}
\item a general interval lower bound for all $c<1$;
\item a new hybrid modular/large-element construction for $1/2<c<1$;
\item an explicit arithmetic lower bound
\[
G_c(n)\ge \left(c-\frac12+\frac{1-c}{q(n)}\right)n+O(1),
\]
where $q(n)$ is the smallest integer $>1$ not dividing $n$.
\end{itemize}

So the Round-1 exact formula for $c\ge 1$ does \emph{not} extend below $1$ by any naive uniform mechanism: already the best explicit constructions for $c<1$ depend on the arithmetic of $n$.

\subsection*{8) COMPLETION ESTIMATE}
COMPLETION: 78\%

\subsection*{9) REFERENCES}
\begin{itemize}
\item T. F. Bloom, \emph{Erd\H{o}s Problem \#361}, \texttt{erdosproblems.com/361} (checked in this round).
\item T. F. Bloom, \emph{Discussion thread for Erd\H{o}s Problem \#361}, \texttt{erdosproblems.com/forum/thread/361} (checked in this round; used only as motivation, not as a proof source).
\item P. Erd\H{o}s and R. L. Graham, \emph{Old and New Problems and Results in Combinatorial Number Theory}, Monographie de L'Enseignement Math\'ematique 28, Geneva, 1980.
\end{itemize}
